\subsection{可标定的工程实现}
\label{sec:engineering}

理想权重要求 $\gamma_{pq}$、$v_q$、$\sigma_{\epsilon,pq}^2$ 在运行时逐像素获得，而单次观测中均不可直接观测；可直接计算的只有 $\Delta_{pq}$、$\hat k_{pq}$、$\rho_{pq}$。本文不为其虚构运行时估计，而是构造与理想似然比同形状的代理，全部自由参数由验证数据离线标定。可靠性因子取
\begin{equation}
  g_{\Delta}(\Delta_{pq})=\exp\!\Big(-\frac{\Delta_{pq}^2}{2\sigma_w^2}\Big),
  \qquad
  g_{\rho}(\rho_{pq})=\rho_{\mathrm{floor}}+(1-\rho_{\mathrm{floor}})\rho_{pq},
  \label{eq:reliability}
\end{equation}
其中 $\sigma_w$ 为几何软门控尺度，标定量级对齐单目局部相对误差（对应 $\sigma_s$）；$\rho_{\mathrm{floor}}\in[0,1)$ 为低分邻居的保底贡献，避免匹配得分偶然偏低时硬性删除仍有先验支持的邻居。精度因子保留显式代理
\begin{equation}
  \hat u_q=\frac{1}{\hat N_q+\hat N_0},
  \qquad
  \hat N_q=\sum_{b\in\Omega_q}\big(H_q(b)-\hat\beta_q\big),
  \label{eq:precision_proxy}
\end{equation}
其中 $\hat N_q$ 为直方图峰显著度，即峰值信号窗口 $\Omega_q$ 内扣除背景估计 $\hat\beta_q$ 后的信号光子数代理，由式~\eqref{eq:v_classical}的反比关系对应理想精度因子 $1/v_q$；$\hat N_0>0$ 为正则化项，对应背景光方差下限，防止低信号邻居权重发散。三个标定参数相互耦合较弱，固定搜索窗口半宽 $K$ 与匹配算子后，在验证集上以深度误差（MAE）为目标逐一维搜索确定并冻结。最终融合权重与归一化系数为
\begin{equation}
  w_{pq}=g_{\Delta}(\Delta_{pq})\,g_{\rho}(\rho_{pq})\,\hat u_q,
  \qquad
  \alpha_{pq}=\frac{w_{pq}}{\sum_{r\in\mathcal{G}_p}w_{pr}},
  \label{eq:weight}
\end{equation}
对对齐后的邻域直方图加权融合
\begin{equation}
  H_p^{\mathrm{fuse}}(b)
  =\sum_{q\in\mathcal{G}_p}\alpha_{pq}\,\widetilde H_{pq}(b),
  \label{eq:fuse}
\end{equation}
再由融合直方图提取回波中心，得到最终深度估计 $\hat\mu_p$。当邻域内信号强度变异有限时，可将 $\hat u_q$ 置为常数，退化为仅可靠性加权。整个流程仅使用单目先验与当前 SPAD 观测即可运行，同时保留了理想贝叶斯模型的核心结构：几何一致性抑制跨表面观测，匹配一致性评价对齐可信度，信号强度压低低光子数邻居的贡献。
